\usepackage[utf8]{inputenc}

%% This is file `sigconf.tex',
%% generated with the docstrip utility.
%%
%% The original source files were:
%%
%% samples.dtx  (with options: `all,proceedings,sigconf')
%% 
%% IMPORTANT NOTICE:
%% 
%% For the copyright see the source file.
%% 
%% Any modified versions of this file must be renamed
%% with new filenames distinct from sigconf.tex.
%% 
%% For distribution of the original source see the terms
%% for copying and modification in the file samples.dtx.
%% 
%% This generated file may be distributed as long as the
%% original source files, as listed above, are part of the
%% same distribution. (The sources need not necessarily be
%% in the same archive or directory.)
%%
%%
%% Commands for TeXCount
%TC:macro \cite [option:text,text]
%TC:macro \citep [option:text,text]
%TC:macro \citet [option:text,text]
%TC:envir table 0 1
%TC:envir table* 0 1
%TC:envir tabular [ignore] word
%TC:envir displaymath 0 word
%TC:envir math 0 word
%TC:envir comment 0 0
%%
%% The first command in your LaTeX source must be the \documentclass
%% command.
%%
%% For submission and review of your manuscript please change the
%% command to \documentclass[manuscript, screen, review]{acmart}.
%%
%% When submitting camera ready or to TAPS, please change the command
%% to \documentclass[sigconf]{acmart} or whichever template is required
%% for your publication.
%%
%%
\documentclass[sigconf]{acmart}
%%
%% \BibTeX command to typeset BibTeX logo in the docs
\AtBeginDocument{%
  \providecommand\BibTeX{{%
    Bib\TeX}}}

%% Rights management information.  This information is sent to you
%% when you complete the rights form.  These commands have SAMPLE
%% values in them; it is your responsibility as an author to replace
%% the commands and values with those provided to you when you
%% complete the rights form.
\setcopyright{acmlicensed}
\copyrightyear{2018}
\acmYear{2018}
\acmDOI{XXXXXXX.XXXXXXX}
%% These commands are for a PROCEEDINGS abstract or paper.
\acmConference[Conference acronym 'XX]{Make sure to enter the correct
  conference title from your rights confirmation email}{June 03--05,
  2018}{Woodstock, NY}
%%
%%  Uncomment \acmBooktitle if the title of the proceedings is different
%%  from ``Proceedings of ...''!
%%
%%\acmBooktitle{Woodstock '18: ACM Symposium on Neural Gaze Detection,
%%  June 03--05, 2018, Woodstock, NY}
\acmISBN{978-1-4503-XXXX-X/2018/06}


%%
%% Submission ID.
%% Use this when submitting an article to a sponsored event. You'll
%% receive a unique submission ID from the organizers
%% of the event, and this ID should be used as the parameter to this command.
%%\acmSubmissionID{123-A56-BU3}

%%
%% For managing citations, it is recommended to use bibliography
%% files in BibTeX format.
%%
%% You can then either use BibTeX with the ACM-Reference-Format style,
%% or BibLaTeX with the acmnumeric or acmauthoryear sytles, that include
%% support for advanced citation of software artefact from the
%% biblatex-software package, also separately available on CTAN.
%%
%% Look at the sample-*-biblatex.tex files for templates showcasing
%% the biblatex styles.
%%
%%
%% The majority of ACM publications use numbered citations and
%% references.  The command \citestyle{authoryear} switches to the
%% "author year" style.
%%
%% If you are preparing content for an event
%% sponsored by ACM SIGGRAPH, you must use the "author year" style of
%% citations and references.
%% Uncommenting
%% the next command will enable that style.
%%\citestyle{acmauthoryear}


%%
%% end of the preamble, start of the body of the document source.
\begin{document}

%%
%% The "title" command has an optional parameter,
%% allowing the author to define a "short title" to be used in page headers.
\title{The Name of the Title Is Hope}

%%
%% The "author" command and its associated commands are used to define
%% the authors and their affiliations.
%% Of note is the shared affiliation of the first two authors, and the
%% "authornote" and "authornotemark" commands
%% used to denote shared contribution to the research.
\author{Ahashan Habib Niloy}
% \authornote{Both authors contributed equally to this research.}
\email{ahashan.niloy@uky.edu}
% \orcid{1234-5678-9012}
% \author{Haibo Wang}
\correspondingauthor
\authornotemark[1]
% \email{wanghaibo@uky.edu}
\affiliation{%
  \institution{University of Kentucky}
  \city{Lexington}
  \state{Kentucky}
  \country{USA}
}

\author{Haibo Wang}
\affiliation{%
  \institution{University of Kentucky}
  \city{Lexington}
  \country{USA}}
\email{wanghaibo@uky.edu}



%%
%% By default, the full list of authors will be used in the page
%% headers. Often, this list is too long, and will overlap
%% other information printed in the page headers. This command allows
%% the author to define a more concise list
%% of authors' names for this purpose.
\renewcommand{\shortauthors}{Niloy et al.}

%%
%% The abstract is a short summary of the work to be presented in the
%% article.
\begin{abstract}
    this is the abstract
 
\end{abstract}

%%
%% The code below is generated by the tool at http://dl.acm.org/ccs.cfm.
%% Please copy and paste the code instead of the example below.
%%
\begin{CCSXML}
<ccs2012>
 <concept>
  <concept_id>00000000.0000000.0000000</concept_id>
  <concept_desc>Do Not Use This Code, Generate the Correct Terms for Your Paper</concept_desc>
  <concept_significance>500</concept_significance>
 </concept>
 <concept>
  <concept_id>00000000.00000000.00000000</concept_id>
  <concept_desc>Do Not Use This Code, Generate the Correct Terms for Your Paper</concept_desc>
  <concept_significance>300</concept_significance>
 </concept>
 <concept>
  <concept_id>00000000.00000000.00000000</concept_id>
  <concept_desc>Do Not Use This Code, Generate the Correct Terms for Your Paper</concept_desc>
  <concept_significance>100</concept_significance>
 </concept>
 <concept>
  <concept_id>00000000.00000000.00000000</concept_id>
  <concept_desc>Do Not Use This Code, Generate the Correct Terms for Your Paper</concept_desc>
  <concept_significance>100</concept_significance>
 </concept>
</ccs2012>
\end{CCSXML}

\ccsdesc[500]{Computing methodologies~Distributed artificial intelligence}
% \ccsdesc[300]{Do Not Use This Code~Generate the Correct Terms for Your Paper}
% \ccsdesc{Do Not Use This Code~Generate the Correct Terms for Your Paper}
% \ccsdesc[100]{Do Not Use This Code~Generate the Correct Terms for Your Paper}

%%
%% Keywords. The author(s) should pick words that accurately describe
%% the work being presented. Separate the keywords with commas.
\keywords{Federated Learning, Multimodality, Communication Efficiency}
%% A "teaser" image appears between the author and affiliation
%% information and the body of the document, and typically spans the
%% page.
% \begin{teaserfigure}
%  \includegraphics[width=\textwidth]{sampleteaser}
%   \caption{Seattle Mariners at Spring Training, 2010.}
%   \Description{Enjoying the baseball game from the third-base
%   seats. Ichiro Suzuki preparing to bat.}
%   \label{fig:teaser}
% \end{teaserfigure}

% \received{20 February 2007}
% \received[revised]{12 March 2009}
% \received[accepted]{5 June 2009}

%%
%% This command processes the author and affiliation and title
%% information and builds the first part of the formatted document.
\maketitle


\section{Introduction}
This is introduction

\section{Related Work}

\subsection{Multimodal Fusion}

\subsection{Compression}


\section{Preliminary}

\subsection{Multimodal Federated Learning}

\subsection{Communication Effiecienet Federated Learning}

\subsection{Selective Federated Learning}

\subsection{Our Proposed Work}



\section{System Design}

\subsection{Overview}

\subsection{Overview}



\section{Evaluation}

\subsection{Experimental Settings}

{\bfseries System Deployment.}
 All experiments were run on a single workstation with an Intel Core i7-13700 CPU (16 cores, 24 threads), 32 GB of RAM and one NVIDIA GeForce RTX 4080 GPU (16 GB), under Ubuntu 24.04 LTS. Federated training is simulated on this machine: the 10 clients and the server share the GPU, and communication is measured analytically as 4 bytes per transmitted 32-bit value. Models are implemented in PyTorch 2.6 (CUDA 12.4, cuDNN 9.1) with Python 3.10; SHAP values are computed with SHAP 0.42 and scikit-learn 1.5 random forests. The FetchSGD baseline uses the authors' CommEfficient codebase on PyTorch 2.5 (CUDA 12.1).

{\bfseries Datasets.}
We evaluate the performance on two multimodal sensing applications (i.e., human activity recognition benchmarks) with three multimodal network designs.
\begin{enumerate}
    \item 
    {\bfseries Smartphone-based Activity Recognition.}
    The UCI Human Activity Recognition dataset (expressed as UCI HAR for short) is used for recognizing daily activities from the inertial sensors of a waist-mounted smartphone. It was collected from 30 volunteers performing 6 activities (walking, walking upstairs, walking downstairs, sitting, standing and laying). Signals are sampled at 50 Hz and segmented into 2.56 s windows with 50\% overlap, and each window is described by 561 time- and frequency-domain features. We split each window's features into two modalities by their source sensor: an accelerometer modality (348-D) and a gyroscope modality (213-D). The dataset includes 7,352 and 2,947 windows for training and test, respectively, following the original subject-disjoint split. 
\end{enumerate}

\begin{enumerate}
    \item 
    {\bfseries Wearable Health Monitoring.}
    The mHealth dataset records body motion and vital signs from 10 subjects performing 12 physical activities. Shimmer2 sensors are placed on the chest, left ankle and right lower arm, sampled at 50 Hz. We discard the null class and segment each contiguous activity run into windows of 128 samples (2.56 s) with 50\% overlap. For every sensor axis and its magnitude channel we extract 14 time-domain statistics (e.g., mean, standard deviation, range, skewness, kurtosis, energy and zero-crossing rate), plus pairwise inter-axis correlations. This yields four modalities: accelerometer (chest, ankle and arm; 177-D), gyroscope (ankle and arm; 118-D), magnetometer (ankle and arm; 118-D), 2-lead ECG (chest; 43-D). Features are z-score normalized using training statistics. To prevent overlapping windows from leaking across the split, we assign whole activity runs to either training or test, which gives 3,681 training and 1,542 test windows.
\end{enumerate}

{\bfseries Models.}
 All models share the same per-modality encoder: a feature extractor with two fully-connected layers ($d_m \rightarrow 512 \rightarrow 512$, with ReLU and dropout of 0.3), followed by a feature refiner with two fully-connected layers ($512 \rightarrow 512 \rightarrow 512$) whose sigmoid output gates the extracted features element-wise. The three designs differ in how they fuse the modalities. Proposed model fuses modalities with a fixed Count Sketch. Each refined 512-D feature vector is hashed into a $2 \times 128$ sketch table, and the tables of all modalities are summed into a 256-D fused representation. An integrator with two fully-connected layers ($256 \rightarrow 512 \rightarrow 512$) and a linear classifier follow. The fusion sketch has no trainable parameters. The model has 2.26M parameters (9.05 MB) on UCI HAR and 3.79M parameters (15.2 MB) on mHealth.

{\bfseries Settings of Statistical Heterogeneity.}
 In the experiments, the training samples are distributed across N = 10 clients following a label-based Dirichlet distribution [cite, e.g., Hsu et al. 2019]. For each class k, we draw a vector p<sub>k</sub> $\sim$ Dir(α·1<sub>N</sub>) that gives the share of class k held by each client, and each training sample of class k is assigned to client j with probability p<sub>k,j</sub>. The concentration parameter α > 0 controls how identical the clients are. As α → ∞, every client receives each class in the same proportion as the global training set (i.e., IID). As α → 0, each class concentrates on a single client, which is a high degree of statistical heterogeneity. We set α = 0.1 to create a highly non-IID setting, and fix the partition across all methods with a common random seed, so every method trains on exactly the same client datasets.

With α = 0.1, the clients are skewed both in label distribution and in dataset size:
- UCI HAR: each client holds between 1 and 5 of the 6 activities, and the most frequent activity makes up 42\% to 100\% of a client's data. Client datasets range from 1 to 1,665 samples (735 on average).
- mHealth: each client holds between 3 and 8 of the 12 activities, with the most frequent activity making up 23\% to 76\% of its data. Client datasets range from 38 to 820 samples (368 on average).

Measured by the KL divergence between each client's label distribution and the global one, the average divergence is 1.24 on UCI HAR and 1.37 on mHealth. The test sets are not partitioned: all methods are evaluated on the same global test set of each dataset (2,947 samples for UCI HAR and 1,542 for mHealth), which covers all classes.

{\bfseries Baselines.}
We measure the effectiveness of our work through a comparison with two approaches.

\begin{enumerate}
    \item 
    {\bfseries Independent Compression}
     is a communication-efficient multimodal FL approach that treats multimodal fusion and gradient compression as two independent steps. Modalities are fused by element-wise summation of their refined features. Each client's full model gradient is then compressed with FetchSGD [Rothchild et al., ICML 2020], i.e., it is projected into a Count Sketch before upload, and the server recovers the top-k coordinates from the aggregated sketch, with momentum and error feedback applied in sketch space. In every round, all clients participate and every modality is uploaded, so the compression is uniform across modalities and ignores how useful each modality or client is to the model.
\end{enumerate}

\begin{enumerate}
    \item 
    {\bfseries  MFedMC}
     [Yuan et al., arXiv:2401.16685] is a state-of-the-art multimodal FL approach that reduces communication through joint modality and client selection. Each client trains a separate unimodal model per modality, then ranks its modalities by a priority score that combines each modality's Shapley value, the size of its model, and how recently it was uploaded. The server selects the clients with the highest local loss. Only the top-ranked modality models of the selected clients are uploaded, as uncompressed model weights, and each modality's global model is aggregated separately. Modalities are fused only at inference, by late fusion of the unimodal predictions. For a fair comparison, we replace MFedMC's original encoders with the same extractor–refiner encoder used by the other methods.
\end{enumerate}

{\bfseries Metrics.}
 We adopt the following metrics to evaluate the performance of our work and the baselines.

\begin{enumerate}
   \item 
  {\bfseries  Test Accuracy}
    reflects the accuracy of the global model trained by each approach, measured as the proportion of the global test set that the model classifies correctly. We evaluate the global model on the same held-out test set after every communication round and record the final accuracy for all approaches. For MFedMC, which keeps a separate model per modality, the prediction is the class with the highest sum of the unimodal log-probabilities.
\end{enumerate}

\begin{enumerate}
   \item 
  {\bfseries   Communication Traffic}
    is the total volume of data exchanged between the clients and the server during training, counting both uploads and downloads, with every transmitted value stored as a 32-bit float.
- Upload: for [Method] and Independent Compression, each participating client uploads one Count Sketch table (rows × columns values). For MFedMC, each selected client uploads the full weights of each selected modality model.
- Download: for the sketch-based methods, a client downloads only the model coordinates that changed since it last participated. For MFedMC, every client downloads all modality models in every round.
- We report the cumulative traffic in megabytes (MB).
\end{enumerate}

\begin{enumerate}
   \item 
  {\bfseries   Communication to Target Accuracy}
    is the cumulative communication traffic an approach consumes before its test accuracy first reaches a target value. For a fair comparison, we set the target to an accuracy that every approach reaches. This metric captures the trade-off between communication and accuracy better than traffic per round, because the approaches differ in how many rounds they run and how much they transmit in each.
\end{enumerate}


{\bfseries Experimental Parameters.}
 By default, all approaches use 10 clients, the extractor–refiner encoder with a 512-D hidden size and a dropout rate of 0.3, and the partition and test sets described above.

\begin{itemize}
\item the [Our work]: each round, one client (10\%) is selected on UCI HAR and two clients (20\%) on mHealth. The selected clients upload the gradient of their top-1 (UCI HAR) or top-2 (mHealth) modalities, computed over their full local dataset. The modality priority weights SHAP value, model size and recency equally (1/3 each). The fusion sketch has 2 rows and 128 columns. The uploaded gradient is compressed into a Count Sketch with 3 rows and 10,000 columns on UCI HAR (5,000 on mHealth), and the server recovers the top k = 20,000 coordinates using momentum 0.9 and error feedback. We run 420 rounds on UCI HAR and 210 on mHealth.
\item Independent Compression: all 10 clients participate in every round and upload a sketch of their full-batch gradient, with 3 rows and 2,500 columns on UCI HAR (5,000 on mHealth). It uses the same k, momentum and error feedback as [Method], and runs 85 rounds on UCI HAR and 42 on mHealth.
\item Learning rate for both sketch-based methods: the server's learning rate increases linearly from 0 to a peak of 0.05 and then decreases linearly back to 0 by the last round. The peak is reached at round 85 (UCI HAR) or 40 (mHealth) for [Method], and at round 20 or 15 for Independent Compression. Weight decay is 5×10⁻⁴ for both.
\item  MFedMC: every client trains each unimodal model for 5 local epochs of SGD, with a fixed learning rate of 0.05 and a batch size of 32. The same fraction of clients as in [Method] is selected (10\% on UCI HAR, 20\% on mHealth), each uploading its top-1 or top-2 modality models, with the same equal priority weights. It runs 7 rounds on UCI HAR and 10 on mHealth.
\item  Seeds: the data partition uses random seed 42 for all approaches; model training uses seed 46 for [Method] and 42 for the baselines.
\end{itemize}



\subsection{Overall Performance}

\subsection{Effect of Non-IID Levels}

\section{Conclusion}




\begin{acks}
This is acknowledgement
\end{acks}



%%
%% The next two lines define the bibliography style to be used, and
%% the bibliography file.
\bibliographystyle{ACM-Reference-Format}
\bibliography{sample-base}



\end{document}
\endinput
%%
%% End of file `sigconf.tex'.
